\section{Results}
\label{sec:results}

\subsection{Main Comparison}

Table~\ref{tab:main} reports held-out perplexity across methods and seeds.
Adaptive reweighting achieves mean PPL 17.651, matching grid-best ($b{=}0.5$, mean 17.646) within seed-to-seed variance: the per-seed gap never exceeds 0.007, less than half the cross-seed range of the uniform baseline (17.686--17.702).
This parity comes at zero validation cost---grid search requires $|\mathcal{B}|$ complete training runs (six in our sweep), whereas adaptive requires one.

The grid-$b$ sweep reveals monotonically increasing PPL from $b{=}0.5$ to $b{=}5.0$ (17.651 to 17.818 on seed 42; full sweep in Appendix~\ref{app:grid}), confirming that exponent selection is consequential.
The optimal $b{=}0.5$ differs from the $b{=}10$ reported by \citet{drayson_2025_detection_collapse} at comparable model scale, underscoring the model-dependence of fixed-exponent selection.

% Claim 1: Dual > ESS-only (3 seeds)
ESS-only (maximum intervention) yields mean PPL 17.723 across three seeds, \emph{worse} than uniform (17.694) by 0.029.
The degradation is consistent: ESS-only exceeds uniform on every individual seed (Table~\ref{tab:main}), ruling out seed-specific artifacts.
Removing $\alpha_{\mathrm{eff}}$ and maximizing $b$ subject only to $\mathrm{ESS} \geq n_{\min}$ leads to systematic over-downweighting, providing direct evidence for the minimum-intervention principle.
Conversely, adaptive (dual) improves over uniform on all three seeds (by 0.041, 0.043, and 0.045 respectively), confirming that the $\alpha_{\mathrm{eff}}$ quality floor prevents the over-intervention that plagues ESS-only.

The Golden Ratio oracle~\citep{he_2025_golden_ratio} achieves PPL 17.654 on seed 42 using known provenance labels.
Adaptive reaches within 0.001 of this upper bound without provenance access, suggesting the dual constraint recovers near-optimal weights from noisy detector scores alone.

The adaptive $b^*$ distribution (median 0.87, IQR $[0.54, 1.51]$, mean 1.40) confirms that most batches need only moderate intervention.
The constraint regime breakdown in \S\ref{sec:algorithm} shows that $\alpha_{\mathrm{eff}}$ prevents over-intervention in the 40\% of reweighted batches where moderate $b$ already meets the quality target.

\begin{table}[t]
\centering
\small
\begin{tabular}{@{}lcccc@{}}
\toprule
\textbf{Method} & \textbf{Seed 42} & \textbf{Seed 123} & \textbf{Seed 456} & \textbf{Mean} \\
\midrule
Uniform               & 17.694 & 17.686 & 17.702 & 17.694 \\
Grid-best ($b{=}0.5$) & \textbf{17.651} & \textbf{17.638} & \textbf{17.650} & \textbf{17.646} \\
ESS-only               & 17.730 & 17.720 & 17.719 & 17.723 \\
Golden Ratio$^\dagger$ & 17.654 & --     & --     & --     \\
\midrule
Adaptive (ours)        & 17.653 & 17.643 & 17.657 & 17.651 \\
\bottomrule
\end{tabular}
\caption{Held-out perplexity (PPL, $\downarrow$).
Grid-best selects the optimal fixed $b$ post-hoc from six candidates.
$^\dagger$Evaluated on seed 42 only.
Best per column in \textbf{bold}.}
\label{tab:main}
\end{table}

\subsection{Sensitivity Analysis}

\paragraph{Quality threshold $\tau$.}
Table~\ref{tab:tau} varies $\tau$ while holding $n_{\min} = 2$ (seed 42).
All values in $\tau \in [0.5, 0.8]$ improve over uniform, with $\tau = 0.8$ achieving the lowest PPL (17.653).
Higher $\tau$ demands greater $\alpha_{\mathrm{eff}}$, forcing the solver to increase $b^*$: mean $\bar{b}^*$ rises monotonically from 0.39 at $\tau{=}0.5$ to 3.08 at $\tau{=}0.95$.
Below 0.8, reweighting is too gentle to fully exploit the detector signal; above, $b^*$ inflates past diminishing returns and PPL degrades.
The range $[0.5, 0.8]$ provides a wide safe operating region, with $\tau = 0.8$ as the recommended default.

\begin{table}[t]
\centering
\small
\begin{tabular}{@{}cccc@{}}
\toprule
$\tau$ & \textbf{PPL} & $\Delta_{\mathrm{unif}}$ & $\bar{b}^*$ \\
\midrule
0.50 & 17.688 & $-$0.006 & 0.39 \\
0.60 & 17.681 & $-$0.013 & 0.59 \\
0.70 & 17.679 & $-$0.015 & 0.89 \\
\textbf{0.80} & \textbf{17.653} & $\mathbf{-0.041}$ & \textbf{1.40} \\
0.90 & 17.712 & $+$0.018 & 2.36 \\
0.95 & 17.723 & $+$0.029 & 3.08 \\
\bottomrule
\end{tabular}
\caption{Sensitivity to $\tau$ (seed 42, $n_{\min}{=}2$, $\rho{=}0.4$).
$\Delta_{\mathrm{unif}}$: PPL change vs.\ uniform baseline (17.694).
$\bar{b}^*$: mean adaptive exponent across batches.}
\label{tab:tau}
\end{table}

\paragraph{Minimum ESS $n_{\min}$.}
Table~\ref{tab:nmin} ablates $n_{\min}$ at batch size $B{=}4$.
Removing the ESS floor ($n_{\min} = 1$) allows $b^*$ to inflate to mean 7.20, degrading PPL to 17.789---0.095 worse than uniform.
At $n_{\min} = 2$ ($= \lfloor B/2 \rfloor$), $b^*$ is constrained to mean 1.40 and PPL reaches 17.653, verified across three seeds.
At $n_{\min} = B = 4$, the ESS constraint is trivially satisfied at $b^* = 0$ for every batch, and PPL reverts to the uniform level (17.695 vs.\ 17.694).
The results exhibit a U-shaped pattern: too-low $n_{\min}$ permits over-downweighting, while $n_{\min}$ at or near $B$ removes the ability to reweight.

The coupling between $n_{\min}$ and $B$ extends beyond the $B{=}4$ setting.
At batch size $B{=}8$ with $n_{\min}{=}2$, the ESS range $[2, 8]$ is wide enough that the stability constraint rarely binds, yielding PPL 17.73---comparable to ESS-only and 0.036 worse than uniform.
This confirms that $n_{\min}$ must scale with $B$: a fixed $n_{\min}$ that works for one batch size can degrade to near-ESS-only behavior at larger $B$.
The rule $n_{\min} = \max(2, \lfloor B/2 \rfloor)$ keeps the ESS constraint active without over-restricting the solver, and we recommend it as a practical default.

\begin{table}[t]
\centering
\small
\begin{tabular}{@{}ccc@{}}
\toprule
$n_{\min}$ & \textbf{PPL} & $\bar{b}^*$ \\
\midrule
1 & 17.789 & 7.20 \\
\textbf{2} & \textbf{17.653} & \textbf{1.40} \\
3$^\ddagger$ & 17.644 & 1.43 \\
4 & 17.695 & 0.00 \\
\bottomrule
\end{tabular}
\caption{Sensitivity to $n_{\min}$ (seed 42, $B{=}4$, $\tau{=}0.8$).
$n_{\min} = 4 = B$ forces $b^* = 0$ on all batches, recovering uniform weighting.
$^\ddagger$From an earlier adaptive configuration; not an independent ablation.}
\label{tab:nmin}
\end{table}
